\documentclass{jsarticle}
\usepackage{amsmath}
\usepackage{amssymb}
\begin{document}
\title{実存主義と思考消去}
\author{Masaaki Yamaguchi}
\date{}
\maketitle
  
   存在と無についての、短めのレポートをこの論文に記す。
   実存主義についての、無と有についての瞑想法を私の思考から導ければいいと
   思って書くことにする。虚無と実存を、虚存と合わせて述べていこうとする。
   まずは、虚存から述べていく。虚存とは、瞑想をするときの周りの場に関係
   している。この場とは、景色に風、匂い、音、そのほかもろもろの瞑想の
   背景となる世界二重内構造を言う。この世界内には、自我が瞑想の実存としての
   陽としてあるのに対して、虚存とは、背景となる陰としての定型をも表している。
   ウィスパードは、この世界二重内構造から投射と関与をもってして、未知の智を
   頭頂葉と側頭葉内前頭葉に思考投入として数式や物理現象、文学を発見する能力が
   ある。この能力は、関与、模倣、投射、接触、破壊、創造、治癒力、その他
   心理学としては、超心理学の分野にこのウィスパードは対象として研究されている。
   自己の心理を他者が関与して読んだり、自己の心理の領域に他者が共鳴力を使って、
   共振として読むのも、速読法を例に取るとわかりやすい。速読法は、自己の思考
   の仕切りを通して、他者の思考の領域と共鳴して、自己と他者が同じ場にいると、
   カタストロフィー現象として、同じ偶然の一致として、同じことを考えていることが
   多い。飼い猫も飼い主の思考が影響して、猫としては思いたくもないが、同じ思考
   になっているらしい。飼い主がだらしないと猫も食べ物がだらしなくなる。
   自然の調和に対して瞑想法ができていると、自己の内部の静けさが音を使わずに、
   瞑想できる。音を使うのも、自分の息に集中するのと同じくらいの支えにもなって
   いる。自己の内部でどのような考え、思考が渦巻いているかを自身で観察して、
   徐々に思考の閾値を下げていき、自然と一体になれば、真逆を功利すると、
   自然と隔離ができるかもしれない。
   一光年三千の天台宗の教えもこの自然との
   共鳴を言い当ててもいる。仏教からの教えは、この世は、生き物が見ている
   その考えの世界への投影に過ぎないと、苫米地博士は、本を通して教えてくれた。
   思っていることがこの世を作っているのですよと、本を通して教えてくれた。
   クレアチニンで思考のデータを消すことも大変ではあるが事実できる人もいる。
   誇大妄想は、この心理の領域に他者が入ると、情報を共有してしまう。
   誇大妄想は、そのまま誇大妄想故に、領域が半端ではない。
   分裂病は、脳に局在化しているかは、このことを言う。
   身体の各部分にも自身の記憶が保存されてもいる。このことも、
   分裂病は、脳に局在化しているかを言い当ててもいる。


\end{document}
